\section*{Capítulo \ref{ch:b2:continuous-reactors} --- Reactores continuos y procesos industriales}

\begin{solution}{exo:b2:continuous-reactors:1}
$\tau = V/Q_v = 2000/0.50 = \qty{4.0e3}{s}$, es decir, \qty{1.1}{h}.
\end{solution}

\begin{solution}{exo:b2:continuous-reactors:2}
$k\tau = 1.2 \times 10^{-3} \times 1800 = 2.16$. Tanque: $X = 2.16/3.16 = 0.68$;
tubo: $X = 1 - \mathrm e^{-2.16} = 0.88$.
\end{solution}

\begin{solution}{exo:b2:continuous-reactors:3}
Tubo: $\tau = \ln 100/0.020 = \qty{230}{s}$. Tanque: $X/(1 - X) = 99 = k\tau$,
$\tau = \qty{4950}{s}$, más de veinte veces más: a \omterm{def:b2:continuous-reactors:conversion}{conversión} alta el
tanque agitado paga caro trabajar a la concentración de salida.
\end{solution}

\begin{solution}{exo:b2:continuous-reactors:4}
Por litro: se liberan $120 \times 0.50 = \qty{60}{kJ}$, que calientan
\qty{4.2}{kJ/K}: $\Delta T_{\text{ad}} = \qty{14}{K}$. Un fallo de la refrigeración
calentaría la mezcla de forma apreciable pero no peligrosa, salvo que el disolvente esté
cerca de su punto de ebullición o que en ese intervalo empiece una segunda reacción.
\end{solution}

\begin{solution}{exo:b2:continuous-reactors:5}
Tanque: $Da = X/(1 - X)^2 = 0.8/0.04 = 20$, $\tau = 20/0.010 = \qty{2000}{s}$.
Tubo: $Da = X/(1 - X) = 4$, $\tau = \qty{400}{s}$. Razón 5, frente a $4/\ln 5 =
2.5$ para el orden 1: cuanto mayor es el orden, más sufre el tanque agitado
por su baja concentración.
\end{solution}

\begin{solution}{exo:b2:continuous-reactors:6}
$k\tau = 5.0$. Tres tanques: $X = 1 - (1 + 5/3)^{-3} = 0.947$. Un tanque: $5/6 =
0.833$. Tubo: $1 - \mathrm e^{-5} = 0.993$.
\end{solution}

\begin{solution}{exo:b2:continuous-reactors:7}
$X = 0.70$ significa $k\tau = 0.70/0.30 = 2.33$. Duplicar el caudal reduce $\tau$ a la mitad:
$k\tau = 1.17$, $X = 0.54$. Para recuperar el 70\,\% hay que duplicar el \omterm{def:b2:continuous-reactors:residence-time}{tiempo de residencia} y,
por tanto, el volumen.
\end{solution}

\begin{solution}{exo:b2:continuous-reactors:8}
Liberado: $80 \times 2.0 \times 1.0 \times 0.90 = \qty{144}{kW}$. Arrastrado por la
corriente de salida, calentada de \qty{330}{K} a \qty{350}{K}: $4.0 \times 1.0 \times 20 =
\qty{80}{kW}$. La camisa debe evacuar la diferencia, \qty{64}{kW}.
\end{solution}

\begin{solution}{exo:b2:continuous-reactors:9}
Dimensiones lineales $\times 10$: volumen $\times 1000$, área de la pared $\times 100$, área
por unidad de volumen $\div 10$. El calor liberado crece con el volumen, y el calor
evacuado a través de la pared, con su área: el reactor grande evacua, por litro,
diez veces menos calor con la misma diferencia de temperatura.
\end{solution}

\begin{solution}{exo:b2:continuous-reactors:10}
Tubo: $\tau = \ln(0.05/0.10)/(0.05 - 0.10) = \qty{13.9}{min}$, $c_B/c_{A,0} =
\frac{0.10}{-0.05}(\mathrm e^{-1.386} - \mathrm e^{-0.693}) = 0.50$. Tanque: $\tau =
1/\sqrt{0.005} = \qty{14.1}{min}$, $c_B/c_{A,0} = 1.414/(2.414 \times 1.707) =
0.34$. En el tanque, A recién llegado y B ya formado están mezclados en el mismo volumen: parte de
B permanece siempre el tiempo suficiente para convertirse en C mientras parte de A apenas acaba de
llegar; en el tubo todas las rebanadas tienen la misma historia.
\end{solution}

\begin{solution}{exo:b2:continuous-reactors:11}
Elevar la temperatura del refrigerante (y de la alimentación) desplaza la recta de evacuación hacia la
derecha. La intersección fría sube lentamente hasta que la recta se vuelve tangente al
codo inferior de la S; más allá, el estado frío desaparece y el tanque salta
a la rama caliente (ignición). Al bajar de nuevo la temperatura, el tanque permanece
en la rama caliente hasta que la recta se vuelve tangente al codo superior, y solo
entonces vuelve a caer (extinción), a una temperatura más baja que la de ignición: un
ciclo de histéresis.
\end{solution}

\begin{solution}{exo:b2:continuous-reactors:12}
Discontinuo: queda el 75\,\% de \qty{2.0}{mol/L}, $1.5 \times 100 = \qty{150}{kJ/L}$,
$\Delta T = 150/4.0 = \qty{37.5}{K}$ por encima de la temperatura de trabajo, lo bastante para
alcanzar un punto de ebullición o una descomposición. Semicontinuo: como mucho
\qty{0.10}{mol/L} sin reaccionar, $\Delta T = 10/4.0 = \qty{2.5}{K}$. La adición lenta
limita la energía almacenada en el reactor.
\end{solution}

\begin{solution}{pb:b2:continuous-reactors:1}
\textbf{1.} $r = k[\text{éster}][\ce{OH-}]$; las alimentaciones iguales y una estequiometría
1 : 1 mantienen iguales las dos concentraciones, $r = kc^2$.
\textbf{2.} $1/c - 1/c_0 = kt$ con $c = 0.10c_0$: $kc_0t = 9$, $t = 9/(0.11
\times 0.10) = \qty{818}{s}$, unos \qty{14}{min}.
\textbf{3.} $t_{1/2} = 1/(kc_0) = \qty{91}{s}$; para el orden 2 la velocidad cae con
el cuadrado de la concentración, así que el tiempo para reducirla a la mitad depende de dónde se
parte.
\textbf{4.} $Da = kc_0\tau = 0.011\,\tau$ ($\tau$ en segundos).
\textbf{5.} $Q_vc_0 - Q_vc - kc^2V = 0$, es decir, $c_0 - c = k\tau c^2$.
\textbf{6.} Con $c = c_0(1 - X)$: $c_0X = k\tau c_0^2(1 - X)^2$.
\textbf{7.} $Da = 0.90/0.10^2 = 90$.
\textbf{8.} $\tau = 90/0.011 = \qty{8.2e3}{s}$, \qty{2.3}{h}.
\textbf{9.} $V = 0.50 \times 8182 = \qty{4.1e3}{L}$, \qty{4.1}{m^3}.
\textbf{10.} La concentración de salida, $0.010\,\unit{mol/L}$: la velocidad es
en todas partes cien veces menor que a la entrada.
\textbf{11.} $Da = X/(1 - X) = 9$, $\tau = \qty{818}{s}$, $V = \qty{409}{L}$.
\textbf{12.} Un volumen diez veces menor para el tubo.
\textbf{13.} $c_{j-1} - c_j = k\tau_1c_j^2$, siendo $\tau_1$ el \omterm{def:b2:continuous-reactors:residence-time}{tiempo de residencia} de
un tanque.
\textbf{14.} $\tau_1 = 850/(3 \times 0.50) = \qty{567}{s}$, $k\tau_1 =
\qty{62.4}{L/mol}$. Resolviendo $62.4c_j^2 + c_j - c_{j-1} = 0$: $c_1 = 0.0328$,
$c_2 = 0.0163$, $c_3 = \qty{0.0100}{mol/L}$: un 90\,\% de \omterm{def:b2:continuous-reactors:conversion}{conversión}.
\textbf{15.} El tubo, o la cascada (diez y cinco veces menores que un solo
tanque). Un solo tanque puede seguir ganando por su temperatura uniforme, su limpieza
y control sencillos, y porque aquí el volumen es barato.
\textbf{16.} $\Delta T_{\text{ad}} = 55\,000 \times 100/(997 \times 4180) =
\qty{1.3}{K}$ (\qty{1.2}{K} al 90\,\%).
\textbf{17.} $55 \times 0.50 \times 0.10 \times 0.90 = \qty{2.5}{kW}$.
\textbf{18.} No: la corriente de salida se lleva el calor con un aumento de aproximadamente
un kelvin; no es posible ningún embalamiento.
\textbf{19.} Veinte veces más concentrada: el aumento adiabático pasa a ser de
\qty{26}{K} y el calor de \qty{50}{kW}, lo que justifica un serpentín de refrigeración; al ser la velocidad $kc_0$
veinte veces mayor, el tanque para un 90\,\% sería veinte veces menor
(\qty{0.20}{m^3}).
\textbf{20.} $0.50 \times 0.10 \times 0.90 = \qty{0.045}{mol/s}$, $\times 86\,400
\times 82.03 = \qty{319}{kg}$ de etanoato de sodio al día.
\textbf{21.} La constante de velocidad crece con la temperatura (Arrhenius), así que el
volumen disminuye; el coste es el calentamiento de la alimentación y, para otras
reacciones, una \omterm{def:b2:continuous-reactors:selectivity}{selectividad} menor y más vapor.
\textbf{22.} El tanque agitado único debe contener $\boldsymbol{V \approx
\qty{4.1}{m^3}}$.
\end{solution}
